% !TeX root = ../../Book1.tex
% 纲要标题：### C.4 有限单群分类概览
% 纲要位置：计划纲要.md，第 877 行。
% BEGIN APPENDIX OUTLINE SYNC
% <!-- 短节：不设小节 -->
% ---
% END APPENDIX OUTLINE SYNC

\section{有限单群分类概览}
\label{b1:sec:C04}

\cref{b1:thm:jordan-holder}说明，有限群的合成因子在同构和排列的意义下确定。因而理解有限单群，是研究有限群结构的一项基本任务。本节概述分类结果及证明中的主要化简；各系列的完整构造和各项子分类的技术证明需要专题文献。

\ProseParagraphBreak
\term[youxian-danqun-fenlei]{有限单群分类}[classification of finite simple groups]断言，每个有限单群都同构于以下名单中的一个群：素数阶循环群，\(n\geq5\) 的交错群 \(A_n\)，有限 Lie 型单群，或二十六个散在单群之一。分类名单见 Aschbacher 的《Sporadic Groups》\cite[Preface, pp. ix--xi]{Aschbacher1994SporadicGroups}。这里将 Tits 群 \({}^2F_4(2)'\) 计入 Lie 型一侧。“Lie 型”和“散在”各自代表具有明确构造的群，不能把这些名称本身当成单性的证明。

\begin{proofsketch}
先排除交换群，它们由\cref{b1:prop:abelian-simple}分类。对其余情形，设有一个阶数最小的不在名单中的单群 \(G\)。任何真子群的简单合成因子都比 \(G\) 小，因此已经在名单中。对非平凡 \(p\)-子群 \(Q\)，正规化子 \(N_G(Q)\) 为真子群；对非中心元 \(x\)，中心化子 \(C_G(x)\) 也为真子群。这使可以对这些局部子群反复使用归纳假设。

奇数阶群定理使非交换 \(G\) 必有偶数阶。随后分析对合的中心化子和二局部子群：低秩情形先作子分类；高秩情形通过局部结构分成有可识别拟单成分和特征二型等情况。这里的拟单成分是中心商为单群的完美次正规子群，它能提供较小的已知简单群及其作用。另一类分析把不同中心化子中满足相容条件的奇数阶子群汇集起来：若汇集出的子群非平凡，其正规化子为真，便严格限制多个局部子群的位置；若汇集部分消失，则转为研究特征二型中的根子群与模。

剩余情形的子分类从这些局部群和它们的交恢复生成元及关系。Lie 型情形由根子群的交换关系识别一个大子群，再证明它就是 \(G\)；特殊局部结构则识别交错群或某个散在群。分类证明还要逐类证明这些识别条件穷尽，不容许最小反例遗漏任何情形。

这里省略奇数阶群定理、各秩的局部子分类、相容子群汇集定理及逐个识别群的长证明；以上说明它们怎样组成排除最小反例的论证，并非给出这些技术定理的完整证明。详细结构及所用论文见 Solomon 的综述\cite[Solvable Signalizer Functor Theorem; Gorenstein--Lyons Trichotomy Theorem]{Solomon2001SimpleGroupsHistory}。各已知群的构造和单性还须分别核验，不能由上述反例化简倒推。
\end{proofsketch}

\ProseParagraphBreak
前两类已经在本卷落实：\cref{b1:prop:abelian-simple}刻画交换单群，\cref{b1:thm:alternating-simple}证明 \(A_n\) 在 \(n\geq5\) 时为非交换单群。Lie 型群把有限域上的矩阵运算及保持某种形式的线性变换发展成若干系列；散在单群则是不属于上述无限系列的有限例外。线性群和 Mathieu 群的具体构造可从 Rotman 的相关两章继续学习\cite[Chapter 8, Some Simple Linear Groups; Chapter 9, Permutations and the Mathieu Groups]{Rotman1995TheoryGroups}。

\ProseParagraphBreak
设 \(\mathbb F_q\) 为有限域。群
\[
\mathrm{SL}_2(\mathbb F_q)
=\left\{\,\begin{pmatrix}a&b\\c&d\end{pmatrix}
\mathrel{}\middle|\mathrel{}a,b,c,d\in\mathbb F_q,\ ad-bc=1\,\right\}
\]
称为二阶特殊线性群，它以矩阵乘法为运算；行列式的乘法性保证乘积仍在其中，通常的二阶逆矩阵公式保证逆元也在其中。其中心恰为满足 \(\lambda^2=1\) 的标量矩阵 \(\lambda I_2\)：这样的矩阵当然在中心；反过来，一个中心元素分别与 \(\begin{psmallmatrix}1&1\\0&1\end{psmallmatrix}\)、\(\begin{psmallmatrix}1&0\\1&1\end{psmallmatrix}\) 交换，直接比较乘积便依次得到 \(c=0,d=a\) 及 \(b=0\)，再由行列式条件得到 \(a^2=1\)。令 \(Z\) 为这个中心，则商群 \(\mathrm{SL}_2(\mathbb F_q)/Z\) 记为 \(\mathrm{PSL}_2(\mathbb F_q)\)，称为二阶\term[sheying-teshu-xianxing-qun]{射影特殊线性群}[projective special linear group]。

\ProseParagraphBreak
这个构造本身不保证单性。例如在 \(\mathbb F_2\) 上，可逆二阶矩阵的第一列有三个选择，第二列有两个不在第一列张成空间内的选择，故共有六个；它们的行列式非零，只能等于一。它们置换 \(\mathbb F_2^2\) 的三个非零向量，且固定这三个向量的矩阵只能是恒等矩阵，因而 \(\mathrm{SL}_2(\mathbb F_2)\cong S_3\)。此时中心平凡，所以 \(\mathrm{PSL}_2(\mathbb F_2)\cong S_3\) 仍不是单群。分类名单中的“有限 Lie 型单群”已经包含对参数与例外的限制，不能把任意有限域上的矩阵商群都不加检查地列入名单。

\ProseParagraphBreak
即使完全知道合成因子，仍不能恢复原群。例如
\[
1<C_2<C_4,\qquad
1<C_2\times\{1\}<C_2\times C_2
\]
都是合成列，因子均为两个 \(C_2\)，但 \(C_4\) 有四阶元素，\(C_2\times C_2\) 的非单位元素全为二阶，所以二者不同构。分类回答了简单因子有哪些，扩张理论还要回答这些因子能够怎样组合。

\ProseParagraphBreak
本卷已证明的 Sylow 定理、可解群性质、交错群单性与 Galois 理论不依赖有限单群分类。第五卷的表示论将提供研究有限群的另一组工具；若进一步阅读有限群结构文献，仍应区分独立的初等证明和以分类为输入的证明。知道分类名单，不能替代核对某一应用实际需要的假设及结论。
% ---
